% problem_id: S001-lemma-Garel-upstairs
% source_id: S001 (Stacks Project)
% source_file: derham.tex
% statement_locator: derham.tex lines 3520--3534
% proof_locator: derham.tex lines 3536--3764
% env: lemma
% id_note: label 在本来源内唯一
% pairing: tight (verified)
% status: complete (statement + author proof verbatim + reason + locators)
%
\begin{lemma}
\label{lemma-Garel-upstairs}
There exists a unique rule that to every locally quasi-finite syntomic
morphism of schemes $f : Y \to X$ assigns $\mathcal{O}_Y$-module maps
$$
c^p_{Y/X} :
\Omega^p_{Y/\mathbf{Z}}
\longrightarrow
f^*\Omega^p_{X/\mathbf{Z}} \otimes_{\mathcal{O}_Y} \det(\NL_{Y/X})
$$
for $p \geq 0$ satisfying (\ref{item-degree-zero}), (\ref{item-multiplicative}),
and (\ref{item-base-change}). In particular, the composition of
$c^p_{Y/X}$ with $f^*\Omega^p_{X/\mathbf{Z}} \to \Omega^p_{Y/\mathbf{Z}}$
is multiplication by $\delta(\NL_{Y/X})$.
\end{lemma}

\begin{proof}
This proof is very similar to the proof of
Discriminants, Proposition \ref{discriminant-proposition-tate-map}
and we suggest the reader look at that proof first.

\medskip\noindent
Let us reformulate the statement. Consider the category
$\mathcal{C}$ whose objects, denoted $Y/X$, are locally quasi-finite syntomic
morphism $f : Y \to X$ of schemes and whose morphisms
$b/a : Y'/X' \to Y/X$ are commutative diagrams
$$
\xymatrix{
Y' \ar[d]_{f'} \ar[r]_b & Y \ar[d]^f \\
X' \ar[r]^a & X
}
$$
which induce an isomorphism of $Y'$ with an open subscheme of
$X' \times_X Y$. The lemma means that for every object
$Y/X$ of $\mathcal{C}$ we have maps $c^p_{Y/X}$, $p \geq 0$
with properties (\ref{item-degree-zero}) and (\ref{item-multiplicative}),
and for every morphism $b/a : Y'/X' \to Y/X$ of $\mathcal{C}$ we have
that $c^p_{Y/X}$ and $c^p_{Y'/X'}$ are compatible as in
Lemma \ref{lemma-base-change-Garel-upstairs}.

\medskip\noindent
Given $Y/X$ in $\mathcal{C}$ and $y \in Y$ we can find
an affine open $V \subset Y$ and $U \subset X$ with $f(V) \subset U$
such that there exists some maps
$$
\Omega^p_{Y/\mathbf{Z}}|_V
\longrightarrow
\left(
f^*\Omega^p_{X/\mathbf{Z}} \otimes_{\mathcal{O}_Y} \det(\NL_{Y/X})
\right)|_V
$$
with properties (\ref{item-degree-zero}) and (\ref{item-multiplicative}).
This follows from picking affine opens as in
Discriminants, Lemma \ref{discriminant-lemma-syntomic-quasi-finite} part (5)
and then using the construction in Remark \ref{remark-local-description}.

\medskip\noindent
Note that the \'etale locus of $f$ is exactly the set of points where
$\delta(\NL_{Y/X})$ does not vanish, see discussion in
Discriminants, Section \ref{discriminant-section-tate-map}.
In particular
$f^*\Omega^p_{X/\mathbf{Z}} \to \Omega^p_{Y/\mathbf{Z}}$
becomes an isomorphism after inverting $\delta(\NL_{Y/X})$.
We conclude that if $\Omega^p_{X/\mathbf{Z}}$ is finite locally free and
the annihilator of $\delta(\NL_{Y/X})$ in $\mathcal{O}_Y$ is zero, then
the local maps constructed in the previous paragraph
are unique and automatically glue!

\medskip\noindent
Let $\mathcal{C}_{nice} \subset \mathcal{C}$ denote the full subcategory
of $Y/X$ such that
\begin{enumerate}
\item $\Omega_{X/\mathbf{Z}}$ is locally free, and
\item the annihilator of $\delta(\NL_{Y/X})$ in $\mathcal{O}_Y$ is zero.
\end{enumerate}
By the remarks in the previous paragraph, we see that for any
object $Y/X$ of $\mathcal{C}_{nice}$ we have unique maps
$c^p_{Y/X}$, $p \geq 0$ satisfying conditions (\ref{item-degree-zero})
and (\ref{item-multiplicative}). If $b/a : Y'/X' \to Y/X$
is a morphism of $\mathcal{C}_{nice}$, then the maps
$c^p_{Y/X}$ and $c^p_{Y'/X'}$ are compatible as in
Lemma \ref{lemma-base-change-Garel-upstairs}: namely,
locally there do exist compatible maps by
Lemma \ref{lemma-existence-base-change-Garel-upstairs}
and the uniqueness just mentioned shows these maps agree
with $c^p_{Y/X}$ and $c^p_{Y'/X'}$.
In other words, we have solved the problem
on the full subcategory $\mathcal{C}_{nice}$. For $Y/X$ in $\mathcal{C}_{nice}$
we continue to denote $c^p_{Y/X}$ the solution we've just found.

\medskip\noindent
In fact, more generally, suppose we have a morphism $b/a : Y'/X' \to Y/X$
of $\mathcal{C}$ such that $Y/X$ is an object of $\mathcal{C}_{nice}$.
The same argument, but this time using the uniqueness in
Lemma \ref{lemma-base-change-Garel-upstairs}, tells us there are
maps $c^p_{Y'/X'}$, $p \geq 0$  satisfying conditions (\ref{item-degree-zero})
and (\ref{item-multiplicative}) compatible with the already constructed
maps $c^p_{Y/X}$. Let us call this the base change
$(b/a)^*c^p_{Y/X}$.

\medskip\noindent
Consider morphisms
$$
Y_1/X_1 \xleftarrow{b_1/a_1} Y/X \xrightarrow{b_2/a_2} Y_2/X_2
$$
in $\mathcal{C}$ such that $Y_1/X_1$ and $Y_2/X_2$ are objects
of $\mathcal{C}_{nice}$.
{\bf Claim.} The two base changes $(b_i/a_i)^*c^p_{Y_i/X_i}$, $i = 1, 2$
are equal. We will first show that the claim implies the lemma
and then we will prove the claim.

\medskip\noindent
Let $d, n \geq 1$ and consider the locally
quasi-finite syntomic morphism $Y_{n, d} \to X_{n, d}$
constructed in Discriminants, Example
\ref{discriminant-example-universal-quasi-finite-syntomic}.
Then $Y_{n, d}$ and $Y_{n, d}$ are irreducible schemes of finite type and
smooth over $\mathbf{Z}$. Namely, $X_{n, d}$ is a spectrum of a
polynomial ring over $\mathbf{Z}$ and $Y_{n, d}$ is an open subscheme
of such. The morphism $Y_{n, d} \to X_{n, d}$ is locally quasi-finite syntomic
and \'etale over a dense open, see Discriminants, Lemma
\ref{discriminant-lemma-universal-quasi-finite-syntomic-etale}.
Thus $\delta(\NL_{Y_{n, d}/X_{n, d}})$ is nonzero: for example we have
the local description of $\delta(\NL_{Y/X})$ in
Discriminants, Remark \ref{discriminant-remark-local-description-delta}
and we have the local description of \'etale morphisms in
Morphisms, Lemma \ref{morphisms-lemma-etale-at-point} part (8).
Now a nonzero section of an invertible module over an irreducible
regular scheme has vanishing annihilator. Thus
$Y_{n, d}/X_{n, d}$ is an object of $\mathcal{C}_{nice}$.

\medskip\noindent
Let $Y/X$ be an arbitrary object of $\mathcal{C}$. Let $y \in Y$.
By Discriminants, Lemma \ref{discriminant-lemma-locally-comes-from-universal}
we can find $n, d \geq 1$ and morphisms
$$
Y/X \leftarrow V/U \xrightarrow{b/a} Y_{n, d}/X_{n, d}
$$
of $\mathcal{C}$ such that $V \subset Y$ and $U \subset X$ are open.
Then we have the base change $c^p_{V/U} = (b/a)^*c^p_{Y_{n, d}/X_{n, d}}$.
The claim guarantees these locally constructed maps $c^p_{V/U}$ glue!
Thus we get a well defined global maps $c^p_{Y/X}$ with properties
(\ref{item-degree-zero}) and (\ref{item-multiplicative}).
Finally, let $b/a : Y'/X' \to Y/X$ be an arbitray morphism of $\mathcal{C}$.
We have the maps $c^p_{Y'/X'}$, $p \geq 0$ and $c^p_{Y/X}$, $p \geq 0$
constructed in this paragraph. To check they are compatible as in
Lemma \ref{lemma-base-change-Garel-upstairs} we may work locally
on $Y'/X'$ and $Y/X$. Thus we may assume that there exists a morphism
$Y/X \to Y_{n, d}/X_{n, d}$. By construction, both the family
$c^p_{Y'/X'}$ and the family $c^p_{X/Y}$ are compatible with
the map towards $Y_{n, d}/X_{n, d}$. It follows easily from this
(and the uniqueness in Lemma \ref{lemma-base-change-Garel-upstairs})
that $c^p_{Y'/X'}$ is compatible with $c^p_{Y/X}$.
Thus it remains to prove the claim.

\medskip\noindent
In the rest of the proof we prove the claim. We may pick a point
$y \in Y$ and prove the maps agree in an open neighbourhood of $y$.
Thus we may replace $Y_1$, $Y_2$ by open neighbourhoods of the
image of $y$ in $Y_1$ and $Y_2$. Therefore we may assume
$Y, X, Y_1, X_1, Y_2, X_2$ are affine, say they are the spectra
of rings $B, A, B_1, A_1, B_2, A_2$. Picture
$$
\xymatrix{
B_1 \ar[r] & B & B_2 \ar[l] \\
A_1 \ar[u] \ar[r] & A \ar[u] & A_2 \ar[l] \ar[u]
}
$$
By assumption the spectrum of $B$ is an affine open of
both the spectrum of $A \otimes_{A_1} B_1$ and $A \otimes_{A_2} B_2$.
Shrinking more we may assume there exist elements
$g_i \in A \otimes_{A_i} B_i$ such that our maps give isomorphisms
$(A \otimes_{A_i} B_i)_{g_i} = B$, see
Properties, Lemma \ref{properties-lemma-standard-open-two-affines}.
Let $x_\alpha$, $y_\beta$ be a sufficiently large collection of
variables such that we may choose surjections
$$
A'_1 = A_1[x_\alpha] \to A
\quad\text{and}\quad
A'_2 = A_2[y_\beta] \to A
$$
of $A_1$ and $A_2$-algebras. Then we can choose lifts
$h_i \in A'_i \otimes_{A_i} B_i$ of $g_i \in A \otimes_{A_i} B_i$
and we consider the diagram
$$
\xymatrix{
(A'_1 \otimes_{A_1} B_1)_{h_1} \ar[r] & B &
(A'_2 \otimes_{A_1} B_2)_{h_2} \ar[l] \\
A'_1 \ar[u] \ar[r] & A \ar[u] & A'_2 \ar[l] \ar[u]
}
$$
By construction the two squares are cocartesian. Next, we consider the
ring map
$$
A' = A'_1 \times_A A'_2 \longrightarrow
B' =
(A'_1 \otimes_{A_1} B_1)_{h_1} \times_B
(A'_2 \otimes_{A_1} B_2)_{h_2}
$$
By
More on Algebra, Lemma \ref{more-algebra-lemma-module-over-fibre-product}
we have $A'_1 \otimes_{A'} B' = (A'_1 \otimes_{A_1} B_1)_{h_1}$
and $A'_2 \otimes_{A'} B' = (A'_2 \otimes_{A_2} B_2)_{h_2}$. In particular
the fibres of the morphism $Y' = \Spec(B') \to \Spec(A') = X'$ are
open subschemes of base changes of the fibres of the maps $Y_i \to X_i$, hence
finite and local complete intersections (see
Discriminants, Lemma \ref{discriminant-lemma-syntomic-quasi-finite}).
By More on Algebra, Lemma
\ref{more-algebra-lemma-properties-algebras-over-fibre-product}
the ring map $A' \to B'$ is flat and of finite presentation.
Thus by Discriminants, Lemma
\ref{discriminant-lemma-syntomic-quasi-finite} part (3)
we conclude that $Y'/X'$ is an object of $\mathcal{C}$.
Consider now the commutative diagram
$$
\xymatrix{
& & Y/X \ar[ld] \ar@/_2pc/[lldd]_{b_1/a_1} \ar[rd] \ar@/^2pc/[rrdd]^{b_2/a_2} \\
& Y'_1/X'_1 \ar[rd] \ar[ld] & & Y'_2/X'_2 \ar[ld] \ar[rd] \\
Y_1/X_1 & & Y'/X' & & Y_2/X_2
}
$$
where $Y'_i/X'_i$ corresponds to $A'_i \to (A'_i \otimes_{A_i} B_i)_{h_i}$.
Note that $Y'_i/X'_i$ is an object of $\mathcal{C}_{nice}$ as it
obtained by taking an open subscheme of the product of an
infinite dimensional affine space with $Y_i/X_i$; small detail omitted.
In particular, the pullback of $c^p_{Y_i/X_i}$ via $(b_i/a_i)$ is the same
as the pullback of $c^p_{Y'_i/X'_i}$ via $Y/X \to Y'_i/X'_i$.
Now we would be done if $Y'/X'$ is an object of $\mathcal{C}_{nice}$,
but this is almost never the case. Namely, then pulling back $c^p_{Y'/X'}$
around the two sides of the square, we would obtain the desired conclusion.
To get around the problem that $Y'/X'$ is not in $\mathcal{C}_{nice}$
we note the arguments above show that, after possibly shrinking all
of the schemes $X, Y, X'_1, Y'_1, X'_2, Y'_2, X', Y'$ we can find some
$n, d \geq 1$, and extend the diagram like so:
$$
\xymatrix{
& Y/X \ar[ld] \ar[rd] \\
Y'_1/X'_1 \ar[rd] & & Y'_2/X'_2 \ar[ld] \\
& Y'/X' \ar[d] \\
& Y_{n, d}/X_{n, d}
}
$$
and then we can use the already given argument by pulling
back from $c^p_{Y_{n, d}/X_{n, d}}$. This finishes the proof.
\end{proof}
